% SPDX-FileCopyrightText: Copyright (c) 2024-2026 Yegor Bugayenko
% SPDX-License-Identifier: MIT

\section{Bytes}\label{sec:bytes}

The center of the taxonomy is the \deff{bytes} object,
  which is an abstraction of a sequence of bytes.
The sequence can either be empty or contain
  a theoretically unlimited number of bytes.
A byte is a unit of information that consists
  of eight adjacent binary digits (bits), each of which consists of a \ff{0}
  or \ff{1}.
In \eolang{} syntax, a sequence of bytes is denoted as either \eohex{---}
  (double dash) for an empty one or as a concatenation of \eohex{xx-},
  where \ff{xx} is a hexadecimal representation of a byte.
For example, the \ff{2A-} is a one-byte sequence,
  where the only byte equals to the decimal number 42.

The \deff{bytes} objects has a few atom attributes for bitwise operations:
  \adeff{and},
  \adeff{or},
  \adeff{not},
  and
  \adeff{right}.
Also, a few object attributes:
  \deff{xor}
  and
  \deff{left}.

The object also has the \adeff{eq} attribute for byte-by-byte comparing
  itself with another sequence of bytes.

The \adeff{size} attribute is the total number of bytes in the sequence.

The \adeff{slice} attribute represents a subsequence inside the current one.

The \adeff{concat} attribute is a new sequence,
  which is a concatenation of two byte sequences:
  the current and the provided one.
